\documentclass[12pt]{article}
\usepackage{amsmath, amssymb}

\title{CSE400 - Fundamentals of Probability in Computing\\
Lecture 20: Linear Transformations and Expectations of Multiple RVs}
\author{Swayam Prajapati, AU2440087}
\date{March 19, 2025}

\begin{document}

\maketitle

\section*{Outline}
\begin{itemize}
    \item Statistical Representation of Random Vectors
    \item Mean Vector, Correlation, and Covariance Matrices
    \item Linear Transformations
    \item Mapping properties under $Y = AX + b$
    \item Multivariate Gaussian Distributions
    \item From 1D/2D cases to N-dimensional Joint PDFs
    \item Special Cases and Properties
    \item Implications of zero correlation and matrix factorization
\end{itemize}
\newpage

\section{Random Vector Representation}

\subsection{Vector Form of Multiple Random Variables}

\textbf{Random Vector}

For random variables $X_1, X_2, \dots, X_N$, we represent them as a vector:

\[
X =
\begin{bmatrix}
X_1 \\
X_2 \\
\vdots \\
X_N
\end{bmatrix}
\]

\begin{itemize}
    \item $X$ is called a \textbf{random vector}.
    \item It provides a compact representation of multiple random variables.
\end{itemize}

\section{Mean Vector}

\subsection{Expectation of a Random Vector}

\textbf{Mean Vector Definition}

The mean vector of $X$ is:

\[
\mu_X = E[X]
\]

\begin{itemize}
    \item Represents the average value of the random vector.
    \item Each element corresponds to the mean of one random variable.
\end{itemize}

\subsection{Example Structure}

\[
\mu_X =
\begin{bmatrix}
E[X_1] \\
E[X_2] \\
\vdots \\
E[X_N]
\end{bmatrix}
\]

\section{Correlation Matrix}

\subsection{Correlation Between Multiple RVs}

\textbf{Correlation Matrix}

\[
R_{XX} = E[X X^T]
\]

where $i, j \leq N$,

\[
R_{XX}(i,j) = E[X_i X_j]
\]

\subsection{Matrix Representation}

\[
R_{XX} =
\begin{bmatrix}
E[X_1 X_1] & E[X_1 X_2] & \cdots & E[X_1 X_N] \\
E[X_2 X_1] & E[X_2 X_2] & \cdots & E[X_2 X_N] \\
\vdots & \vdots & \ddots & \vdots \\
E[X_N X_1] & E[X_N X_2] & \cdots & E[X_N X_N]
\end{bmatrix}
\]

\begin{itemize}
    \item $R_{XX}$ is a symmetric matrix.
\end{itemize}

\section{Covariance Matrix}

\subsection{Mathematical Definition}

For a random vector $X$, the covariance matrix $C_{XX}$ is:

\[
C_{XX} = E[(X - \mu_X)(X - \mu_X)^T]
\]

\subsection{Matrix Form in Practice}

\[
C_{XX} =
\begin{bmatrix}
\mathrm{Var}(X_1) & \mathrm{Cov}(X_1, X_2) & \cdots \\
\mathrm{Cov}(X_2, X_1) & \mathrm{Var}(X_2) & \cdots \\
\vdots & \vdots & \ddots
\end{bmatrix}
\]

\begin{itemize}
    \item Symmetry: $\mathrm{Cov}(X_i, X_j) = \mathrm{Cov}(X_j, X_i)$
\end{itemize}

\section{Linear Transformation of Random Vectors}

\subsection{The Linear Model}

Consider the transformation:

\[
Y = AX + b
\]

\begin{itemize}
    \item $A$: Transformation matrix (scales and rotates the data)
    \item $b$: Constant vector (shifts/translates the center)
    \item $X$: Input random vector
\end{itemize}

For $N$ dimensional vectors:

\[
Y_1 = a_{11}X_1 + a_{12}X_2 + \cdots + a_{1N}X_N + b_1
\]
\[
Y_2 = a_{21}X_1 + a_{22}X_2 + \cdots + a_{2N}X_N + b_2
\]
\[
\vdots
\]
\[
Y_N = a_{N1}X_1 + a_{N2}X_2 + \cdots + a_{NN}X_N + b_N
\]

\section{Mean of Transformed Vector}

\subsection{Mean Transformation}

For a linear transformation $Y = AX + b$:

\[
\mu_Y = E[Y]
\]

\subsection{Derivation}

Substitute $Y = AX + b$ into the expectation:

\[
\mu_Y = E[AX + b]
\]

Apply linearity:

\[
E[AX] + E[b]
\]

\subsection{Result}

\[
\mu_Y = A \mu_X + b
\]

\section{Correlation Matrix Transformation}

\subsection{Definition}

\[
R_{YY} = E[Y Y^T] = E[(AX + b)(AX + b)^T]
\]

\subsection{Expansion of Terms}

\[
R_{YY} = E[AXX^T A^T + AXb^T + bX^T A^T + bb^T]
\]

\subsection{Applying Expectation}

\[
R_{YY} = A E[XX^T] A^T + A E[X] b^T + b E[X^T] A^T + bb^T
\]

\subsection{Final Result}

\[
R_{YY} = A R_{XX} A^T + A \mu_X b^T + b \mu_X^T A^T + bb^T
\]

\section{Covariance Matrix Transformation}

\subsection{Key Observation}

The deviation from the mean is independent of the shift $b$:

\[
Y - \mu_Y = (AX + b) - (A\mu_X + b) = A(X - \mu_X)
\]

\end{document}